\scanpage{103}
\mspage{51}
On a évidemment $\mathrm{a)}\Rightarrow\mathrm{b)}\longrightarrow\mathrm{c)}\unsure{=|\Rightarrow}\mathrm{d)}$, et $d\Rightarrow a$ devrait être plus ou moins connu, ou du moins se prouver à l'aide de choses bien connues. Ces conditions impliquent e, mais je n'ai pas d'idée pour le moment comment voir que e) implique bonne réduction --- faudrait une compréhension (\unsure{fut-elle} heuristique) plus approfondie que celle que j'ai pour le moment.
\margin{Ça vaudrait la peine de regarder une \unsure{version} de ceci où $X_\eta=\bar U_\eta$ a bonne réduction --- mais la multisection $\overline{X_\eta-U_\eta}$ de $X$ sur $S$ n'est pas étale. Le cas du genre $0$ devrait être particulièrement aisé à regarder\ldots}

J'ai donc là un principe général (que je pourrais à l'occasion formuler avec la généralité qui lui convient) pour construire des actions \unsure{extérieures|\deleted{extérieures}} \inserted{Individuelles} de groupes $\mathcal E_k$ ($k$ \deleted{ext.} de type fini de $\mathbf Q$) sur des $\pi$ à \unsure{bouts}, qui \inserted{(sans doute)} ne soient géométriques, par \deleted{des} la considération de courbes de type $g,\nu$ ``se réduisant mal''. Mais je \unsure{n'arrive|reviens}
\margin{102}
